%%%BLOCK id=170-01 ch=18 sec=力学类实验·总览 kind=summary alt=01,02,03,06,07 cont=none y=0.03-0.19
\textbf{力学类实验}
\begin{itemize}
\item 弹簧
\item 力平行四边形
\item 测$a$.（测$a$时不用测$m$，也不平衡$f$）\quad 电火花通电\quad 电磁用永磁体\quad\textit{（该行上方小字：\unclear{指出操作有误}$\downarrow$\quad $220\unit{V}$\quad $6\sim8\unit{V}$）}
\item 探究$aFm$关系
\item 测$M$
\item 能量
\item 动量
\end{itemize}
%%%END
%%%BLOCK id=172-04 ch=18 sec=气垫导轨·使用注意 kind=note exp=1 alt=03,07 cont=none y=0.86-0.91
气垫导轨：先拿滑块，后关电源，防止摩擦损 % CHECK: 句末「损」字下笔画涂改，疑为「损坏」未写完
%%%END
%%%BLOCK id=182-03 ch=18 sec=气垫导轨·判断导轨水平 kind=method exp=1 alt=03,06 cont=none y=0.425-0.485
判断气垫导轨水平的方法：接通电源，将滑块静置于气垫导轨上，①不放钩码时，若滑块保持静止②，或轻推滑块，滑块做匀速直线运动则说明导轨水平。 % 原稿“钩码”二字写得像“约码”；圈号①②为小圈号，位置如此
%%%END
%%%BLOCK id=183-03 ch=18 sec=数据处理·描点作图 kind=method exp=1 alt=none cont=none y=0.40-0.47
坐标描点时 按等差数列找

选点靠\unclear{…}找最远点

0.3\quad 0.6\quad 0.9\quad 1.2\quad $\cdots$\quad 2.1

0.01\quad 0.02\quad 0.03\quad $\cdots$\quad 0.07
%%%END
%%%BLOCK id=193-10 ch=18 sec=游标卡尺读数 kind=note alt=none cont=none y=0.86-0.93
游标卡尺
%FIG p193_e 0.12 0.862 0.265 0.912
\notefig[0.25]{figures/p193_e.png}

\underline{不是从0开始}
%%%END
%%%BLOCK id=219-04 ch=18 sec=实验方法·力学实验 kind=note exp=1 alt=none cont=none y=0.565-0.61
力学实验当成力学计算题
%%%END
%%%BLOCK id=221-11 ch=18 sec=实验·计数计时 kind=note exp=1 alt=08 cont=none y=0.96-1.00
1到16次 \unclear{算}时 15 除以15 % 原稿："15"上方有三点及横线
%%%END
%%%BLOCK id=234-01 ch=18 sec=实验·答题方法 kind=method exp=1 alt=none cont=none y=0.07-0.20
\textbf{实验}\par
% “忆课内实验”为写在“联想”上方的小字
实验目的\ 首句精读 \quad 实验原理\ 通读全篇 \quad \unclear{忆}课内实验\ 联想 \quad 关注小括号（\unclear{微}数字）代号or字母\par
器材\ 步骤\ 数据处理\ 误差分析
%%%END
%%%BLOCK id=245-06 ch=18 sec=游标卡尺·读数方法 kind=method alt=none cont=none y=0.86-1.00
\unclear{游}标卡尺\quad 减法：上尺对齐\unclear{读数}$-$格数$(1-\text{分度值})$ % 原稿此行有几处浅色涂抹痕迹，字仍可见；首字被页边截断

\unclear{分度}值$=\dfrac{1}{\text{格数}}$\quad 加法：下尺对整数$+$格数$\times$分度值 % 第二行行首被页边截断

与格数无关 % 其右侧有一处极小的涂写（疑似“10”或“1/0”），无法辨认
%%%END
%%%BLOCK id=249-05 ch=18 sec=数据处理·图像规范 kind=note alt=none cont=none y=0.86-0.92
图像 $\tan\theta$ 值不等于切线斜率\quad 分度值大小不规范 \\
物理量$+$单位
%%%END
%%%BLOCK id=273-06 ch=18 sec=读数·易错 kind=note alt=none cont=none y=0.72-0.85
\textbf{读数坑}

%FIG p273_d 0.03 0.745 0.25 0.815
\notefig[0.3]{figures/p273_d.png}

$\Delta x=x_2-2x_1$
%%%END
%%%BLOCK id=273-07 ch=18 sec=游标卡尺·读数 kind=note alt=none cont=none y=0.86-0.89
游标卡尺副尺几个格就是几分度 % 原稿“副尺”为插入的小字，写在“尺”“几”之间的上方
%%%END
%%%BLOCK id=295-03 ch=18 sec=数据处理·选点 kind=note exp=1 alt=none cont=none y=0.53-0.56
实验\underline{题选点}必选最远点。
%%%END
